% SCSE / HRAI-GO | Especificação técnica ESP-SCSE-01 (ABNT NBR 10719:2015)
% Compilação: scripts/build-relatorios.sh docs/scse/SCSE_Shadow_Council_Simulation_Spec.tex
% (XeLaTeX + BibTeX via latexmk). Regras: docs/PADRAO_RELATORIOS.md
\documentclass[12pt,a4paper,oneside,english,brazil]{abntex2}
\usepackage{hraigo-abnt}

% ------------------------------------------------------------ Pacotes e macros próprios deste documento
\usepackage{amsmath}
\usepackage{listings}
\usepackage{seqsplit,needspace}
\renewcommand{\arraystretch}{1.22}
\setlength{\LTcapwidth}{\textwidth}  % longtable captions as wide as the text
\setlist[itemize]{leftmargin=16pt,topsep=3pt,itemsep=2pt}
\setlist[enumerate]{leftmargin=19pt,topsep=3pt,itemsep=2pt}
% Code listings: single spacing, house palette.
\lstdefinestyle{code}{basicstyle=\ttfamily\scriptsize\SingleSpacing,columns=fullflexible,breaklines=true,breakatwhitespace=false,keepspaces=true,showstringspaces=false,frame=single,rulecolor=\color{tinta!23},backgroundcolor=\color{cinza!6},framesep=7pt,aboveskip=8pt,belowskip=8pt,upquote=true}
\lstset{style=code}
% \term{...}: canonical identifier in monospace, breakable at any character.
\newcommand{\term}[1]{{\ttfamily\small\seqsplit{#1}}}
% \mvp: bold acronym of the minimum viable product.
\newcommand{\mvp}{\textbf{MVP}}

% ------------------------------------------------------------ Identificação (NBR 10719, 4.2.1.1.1)
\titulo{Shadow Council Simulation Environment}
\subtitulo{ambiente deliberativo multiagente inspirado no conceito do MiroFish, para avaliar regulações e soluções perante regimes jurídicos específicos}
\codigo{ESP-SCSE-01}
\tipotrabalho{Especificação técnica}
\versao{1.1}
\data{2026}
\palavraschave{Shadow Council. Simulação multiagente. Inteligência regulatória. Conformidade regulatória. Proveniência de evidências.}
\hypersetup{pdftitle={SCSE: Shadow Council Simulation Environment},pdfsubject={Especificação técnica ESP-SCSE-01 do HRAI-GO}}

\begin{document}
\selectlanguage{brazil}

% ============================================================ ELEMENTOS PRÉ-TEXTUAIS
\pretextual
\imprimirfolhaderostorelatorio

\begin{resumorelatorio}
Esta especificação técnica descreve o Shadow Council Simulation Environment (SCSE), módulo independente do
HRAI-GO (Human-centred Regulatory AI, Governance Observatory). O SCSE é um ambiente computacional isolado que,
diante de uma semente de realidade (regulação, proposta normativa, briefing ou especificação), constrói um
estado de mundo versionado, seleciona perspectivas sintéticas documentadas, executa um processo de
contestação multiagente e emite relatórios reproduzíveis, com evidências, divergências e limitações. A
questão de produto é identificar que obrigações, riscos, lacunas, tensões e perguntas críticas emergem quando
uma regulação ou solução é examinada por um conselho deliberativo sintético contextualizado por uma
jurisdição. O documento separa a análise de adequação de política da verificação de conformidade legal e
especifica a composição estratificada do conselho, o modelo canônico de estado, o protocolo deliberativo, o
motor jurídico, os contratos de dados, a interface de programação, os requisitos verificáveis, o método de
avaliação e a ordem de implementação. Ficam fora de escopo a certificação jurídica automática, a consulta
pública substitutiva, a previsão estatística da opinião social, agentes que assumam a identidade de pessoas
reais e simulações que alterem sistemas reais. O documento utiliza apenas interfaces contratuais com o
HRAI-GO e com catálogos externos, e distingue exemplos sintéticos e requisitos propostos de informações sobre
ferramentas e normas existentes. O conselho sintético não substitui pessoas, consultas públicas ou parecer
jurídico. Trata-se de especificação arquitetural e funcional, não de uma implementação executada.
\end{resumorelatorio}

\pdfbookmark[0]{\listtablename}{lot}
\listoftables*
\cleardoublepage

\pdfbookmark[0]{\contentsname}{toc}
\tableofcontents*
\cleardoublepage

% ============================================================ ELEMENTOS TEXTUAIS
\textual

\chapter{Introdução: visão geral e limite do módulo}
\noindent\textit{Nome da plataforma.} HRAI-GO significa \textit{Human-centred Regulatory AI, Governance Observatory} (plataforma de inteligência regulatória humano--IA); ver \texttt{docs/NOMENCLATURA.md}.

Este documento é a especificação funcional, semântica, multiagente e operacional de um módulo independente
do HRAI-GO (Human-centred Regulatory AI, Governance Observatory). O texto está em português do Brasil; os
identificadores canônicos ficam em inglês.

\destaque{Documento independente: utiliza apenas interfaces contratuais com o HRAI-GO e com catálogos externos. Exemplos sintéticos e requisitos propostos são distinguidos de informações sobre ferramentas e normas existentes. O conselho sintético não substitui pessoas, consultas públicas ou parecer jurídico.}

\destaque{\textbf{Definição.} \term{ShadowCouncilSimulationEnvironment} (\term{SCSE}) é um ambiente computacional isolado que, diante de uma \emph{semente de realidade} (regulação, proposta normativa, briefing ou especificação), constrói um estado de mundo versionado, seleciona perspectivas sintéticas documentadas, executa um processo de contestação multiagente e emite relatórios reproduzíveis, com evidências, divergências e limitações.}

\noindent\textbf{Questão de produto:} Que obrigações, riscos, lacunas, tensões e perguntas críticas emergem quando uma regulação ou solução é examinada por um conselho deliberativo sintético contextualizado por uma jurisdição?

A Figura~\ref{fig:fluxo} resume o fluxo do SCSE, das sementes de realidade aos relatórios com opiniões minoritárias.

\begin{figure}[htb]
\caption{Fluxo geral do SCSE: das sementes de realidade ao relatório com opiniões minoritárias}\label{fig:fluxo}
\centering
\begin{tikzpicture}[x=1cm,y=1cm,>=Latex]
\fill[tinta!5,rounded corners=7pt] (0,0) rectangle (15.8,5.9);
\tikzset{source/.style={draw=tinta!65,fill=white,rounded corners=4pt,minimum width=6.5cm,minimum height=.95cm,align=center,font=\sffamily\bfseries\small},chamber/.style={draw=lagoa!75,fill=white,rounded corners=4pt,minimum width=6.5cm,minimum height=1cm,align=center,font=\sffamily\bfseries\small}}
\node[source] (i) at (7.9,5.1) {Reality Seeds + Regulatory Snapshot};
\node[source] (w) at (7.9,3.7) {World Graph + Stakeholder Sampling Frame};
\node[chamber] (p) at (3.9,2.18) {Policy Deliberation Chamber};
\node[chamber] (c) at (11.9,2.18) {Solution Compliance Chamber};
\node[source] (o) at (7.9,.63) {Evidence-Aware Report + Minority Opinions};
\draw[->,thick,tinta!65] (i)--(w);
\draw[->,thick,tinta!65] (w.south)--(p.north);
\draw[->,thick,tinta!65] (w.south)--(c.north);
\draw[->,thick,lagoa!80] (p.south)--(o.north);
\draw[->,thick,lagoa!80] (c.south)--(o.north);
\end{tikzpicture}
\fonte{Elaborado pelos autores (2026).}
\end{figure}

\section{Objetivos observáveis}
\begin{enumerate}
\item Aceitar um \term{RegulatoryContext} ou um \term{SolutionArtifact} (briefing, spec, arquitetura), com país e data de referência.
\item Separar análise de \emph{adequação de política} de \emph{verificação de conformidade legal}.
\item Compor um \term{ShadowCouncil} estratificado por interesses, papéis e perspectivas documentadas.
\item Produzir argumentação contraditória, testes de hipóteses, relatórios minoritários e trilhas de evidência.
\item Permitir simular contrafactuais com alterações controladas (\term{ScenarioIntervention}) e comparar execuções.
\item Preservar \term{RunManifest}, estado inicial, semente, versões de modelos, documentos, prompts, eventos e resultados.
\end{enumerate}
\section{Fora de escopo (na versão inicial)}
\begin{itemize}
\item Certificação jurídica automática, consulta pública substitutiva ou previsão estatística da opinião social.
\item Agentes que assumam identidade ou mandato de autoridades ou pessoas reais sem autorização.
\item Simulação que altere sistemas reais, publique textos em redes sociais ou execute mudanças em produção.
\item Conclusão quantitativa de causalidade regulatória sem modelo empiricamente calibrado.
\end{itemize}

\chapter{O que reaproveitar do conceito MiroFish}
A documentação do \href{https://github.com/666ghj/MiroFish}{projeto original MiroFish} \cite{mirofish} apresenta cinco passos: \emph{graph building}, \emph{environment setup}, \emph{simulation}, \emph{report generation} e \emph{deep interaction}. Seu simulador social usa o \href{https://github.com/camel-ai/oasis}{OASIS}, desenvolvido pela CAMEL-AI \cite{camel-oasis}. O SCSE considera essa referência arquitetural, sem afirmar que a base existente já implemente raciocínio jurídico. A Tabela~\ref{tab:mirofish} mostra como cada conceito original é adaptado.

\begin{table}[!htb]
\caption{Adaptação dos conceitos do MiroFish no SCSE}\label{tab:mirofish}
\centering
\begin{tabularx}{\textwidth}{@{}p{.23\textwidth}L L@{}}
\toprule\textbf{Conceito original} & \textbf{Adaptação SCSE} & \textbf{Critério}\\\midrule
Reality seeds & \term{EvidenceBundle} + textos jurídicos + especificação & Proveniência e estado temporal obrigatórios\\
Knowledge graph & \term{WorldSnapshot} com normas, stakeholders, aplicação e argumentos & Ontologia canônica + dados de sessão\\
Personas e memória & \term{PerspectiveAgent} + \term{PerspectiveCharter} + memória isolada & Cobertura documental, não imitação humana\\
Interações sociais & \term{DeliberationChamber} com objeção, prova, emenda e revisão & Rodadas estruturadas e auditáveis\\
Relatório e chat & \term{CouncilReport} + rastreio de achados + interrogatório de agentes & Sem falsa inferência de voto popular\\\bottomrule
\end{tabularx}
\fonte{Elaborado pelos autores (2026), com base em \citeonline{mirofish}.}
\end{table}

\noindent\textbf{Decisão de implementação.} Para o produto mínimo viável (\mvp, \textit{minimum viable product}): construir motor deliberativo próprio que exponha \term{SimulationBackendAdapter}; o OASIS pode ser um \emph{adapter} para experimentos de difusão e interação, mas não será dependência obrigatória para conformidade. Se reutilizar código MiroFish, respeitar sua licença GNU Affero General Public License (AGPL) versão 3.0 e analisar obrigações de distribuição; a inspiração conceitual não exige incorporar código.

\chapter{Modos de simulação e contrato semântico}
\section{Mode A: Regulatory Policy Simulation}
\noindent\textbf{Entrada:} nação, data, versão de norma ou proposta, objetivos de política, fontes de contexto nacional, pergunta deliberativa e desenho amostral.\\
\textbf{Pergunta-tipo:} ``Quais grupos, objetivos e capacidades institucionais seriam afetados por um regime brasileiro de avaliação de impacto de IA?''\\
\textbf{Saída:} \term{RegulatoryTension}, \term{RegulatoryGapCandidate}, argumentos a favor/contra, \term{PolicyAlternative}, condições de viabilidade e evidências faltantes.\\
\textbf{Regra epistêmica:} opiniões dos agentes não implicam validade jurídica ou legitimidade democrática.

\section{Mode B: Solution Compliance Simulation}
\noindent\textbf{Entrada:} briefing/spec de solução, país e setor, data, marco regulatório selecionado ou pacote aplicável, arquitetura, fluxos de dados, atores, operações, riscos e evidências disponíveis.\\
\textbf{Pergunta-tipo:} ``O que falta demonstrar para que um sistema de triagem de candidatos esteja conforme com as exigências aplicáveis no Brasil?''\\
\textbf{Saída:} \term{ApplicabilityMatrix}, \term{ComplianceFinding}, \term{EvidenceGap}, requisitos em aberto, questões jurídicas, remediações propostas e argumentos dos stakeholders.\\
\textbf{Regra epistêmica:} \term{LegalAssessmentEngine} determina aplicabilidade e estado probatório; \term{CouncilSimulationEngine} sugere objeções e alternativas sem sobrescrever os achados verificáveis.

\section{Tabela de diferenças obrigatórias}
A Tabela~\ref{tab:modos} resume as diferenças obrigatórias entre os dois modos.

\begin{table}[!htb]
\caption{Diferenças obrigatórias entre o modo de política e o modo de solução}\label{tab:modos}
\centering
\begin{tabularx}{\textwidth}{@{}p{.25\textwidth}LL@{}}
\toprule & \textbf{Policy mode} & \textbf{Solution mode}\\\midrule
Objeto & qualidade da política & aplicação/projeto de IA\\
Referência & contexto nacional + proposta/norma & norma vigente na data + objeto técnico\\
Resultado & deliberação, tensões e alternativas & matriz normativa + avaliação de evidência\\
Conselho & contrasta efeitos sociais/institucionais & contesta pressupostos de conformidade e impacto\\
Estado final & \term{PolicyScenarioReport} & \term{SolutionComplianceReport}\\\bottomrule
\end{tabularx}
\fonte{Elaborado pelos autores (2026).}
\end{table}

\chapter{Arquitetura de referência independente}
A Figura~\ref{fig:arquitetura} apresenta a arquitetura proposta.

\begin{figure}[htb]
\caption{Arquitetura proposta. Os mecanismos jurídico e deliberativo compartilham evidências, mas possuem saídas e autoridade distintas}\label{fig:arquitetura}
\centering
\begin{tikzpicture}[x=1cm,y=1cm,>=Latex]
\tikzset{b/.style={draw=tinta!68,fill=cinza!6,rounded corners=2pt,minimum height=.9cm,text width=4cm,align=center,font=\sffamily\footnotesize},wide/.style={draw=lagoa!65,fill=lagoa!8,rounded corners=3pt,minimum height=.84cm,text width=13.7cm,align=center,font=\sffamily\footnotesize\bfseries}}
\node[wide] (ui) at (7.5,8.7) {SCSE Studio / API / CLI / External Agent Gateway};
\node[b] (ing) at (2.4,7) {Ingestion + Citation Resolver};
\node[b] (kg) at (7.5,7) {World Builder + Graph Snapshot};
\node[b] (sampling) at (12.6,7) {Stakeholder Sampling + Persona};
\node[wide] (orch) at (7.5,5.55) {Scenario Orchestrator + Event Ledger + Run Manifest};
\node[b] (legal) at (2.4,3.95) {Legal Assessment Engine\\OWL / SHACL / Rules};
\node[b] (sim) at (7.5,3.95) {Council Simulation Engine\\Rounds + Memory};
\node[b] (verify) at (12.6,3.95) {Evidence / Critic / Safety};
\node[wide] (reports) at (7.5,2.35) {Reports + Dissent + Sensitivity + Provenance};
\node[wide] (stores) at (7.5,.8) {Object Store | Graph Store | Event Store | Vector Index};
\draw[->,thick,tinta!70] (ui)--(kg);
\draw[->,thick,tinta!70] (ing.south)--(orch.north -| ing);
\draw[->,thick,tinta!70] (kg)--(orch);
\draw[->,thick,tinta!70] (sampling.south)--(orch.north -| sampling);
\draw[->,thick,tinta!70] (orch.south)--(sim.north);
\draw[->,thick,tinta!70] (orch.south -| legal)--(legal.north);
\draw[->,thick,tinta!70] (orch.south -| verify)--(verify.north);
\draw[->,thick,lagoa] (sim.south)--(reports.north);
\draw[->,thick,lagoa] (legal.south)--(reports.north -| legal);
\draw[->,thick,lagoa] (verify.south)--(reports.north -| verify);
\draw[<->,thick,lagoa] (reports)--(stores);
\end{tikzpicture}
\fonte{Elaborado pelos autores (2026).}
\end{figure}

\section{Componentes e isolamento}
A Tabela~\ref{tab:componentes} lista os componentes, com o contrato e a responsabilidade de cada um.

\begin{longtable}{@{}p{.34\textwidth}p{.58\textwidth}@{}}
\caption{Componentes do SCSE, contratos e responsabilidades}\label{tab:componentes}\\
\toprule\textbf{Componente} & \textbf{Contrato e responsabilidade}\\\midrule\endfirsthead
\toprule\textbf{Componente} & \textbf{Contrato e responsabilidade}\\\midrule\endhead
\term{SeedIngestionService} & Normalizar PDF, Markdown (MD), TXT, JavaScript Object Notation (JSON), links oficiais e specs; preservar trecho, hash, identificador uniforme de recurso (URI) e estrutura do documento.\\
\term{LegalSnapshotResolver} & Resolver \term{validTime}/\term{transactionTime}, jurisdição, status de vigência, hierarquia e versão; falha fechada para fonte essencial não verificada.\\
\term{WorldModelBuilder} & Gerar \term{WorldSnapshot} congelado a partir de normas, contexto nacional e entidade avaliada, com entidades tipadas e claims rastreáveis.\\
\term{CouncilComposer} & Definir estratos, cobertura, fontes, assentos, agentes e limites de atuação.\\
\term{SimulationOrchestrator} & Gerir rounds, seeds, branches, intervenções, custos, checkpoints e reexecuções.\\
\term{LegalAssessmentEngine} & Determinar obrigações aplicáveis e produzir achados sob status controlados, usando regras formais e validação.\\
\term{DeliberationEngine} & Coordenar discussão, objeções, revisão, negociação, dissenso e análise contrafactual.\\
\term{EvidenceVerifier} & Verificar citações, cobertura dos trechos, validade temporal e contradições materiais.\\
\term{ReportCompiler} & Gerar JSON, MD e PDF com trilhas de proveniência e indicadores de incerteza.\\
\term{ExternalAgentGateway} & Receber agentes de terceiros por capacidades e permissões, com isolamento de segurança.\\\bottomrule
\end{longtable}
\fonte{Elaborado pelos autores (2026).}

\chapter{Modelo canônico de estado e ontologia}
\section{Entidades essenciais}
A Tabela~\ref{tab:entidades} define a semântica das entidades essenciais do modelo.

\begin{longtable}{@{}p{.37\textwidth}p{.55\textwidth}@{}}
\caption{Entidades essenciais do modelo canônico de estado}\label{tab:entidades}\\
\toprule\textbf{Classe} & \textbf{Semântica}\\\midrule\endfirsthead
\toprule\textbf{Classe} & \textbf{Semântica}\\\midrule\endhead
\term{SimulationProject} & Experimento, pergunta, proprietário, escopo e controle de acesso.\\
\term{ScenarioDefinition} & Modo, jurisdição, as-of-date, objetivos e hipóteses.\\
\term{RegulatorySnapshot} & Pacote normativo em estado juridicamente contextualizado.\\
\term{NationalContextSnapshot} & Recorte temporal de indicadores institucionais e sociais.\\
\term{SolutionArtifact} & Briefing, especificação, arquitetura ou ficha da aplicação.\\
\term{EvidenceRecord} & Unidade citável com fonte, trecho, data, hash, tipo e confiabilidade.\\
\term{StakeholderStratum} & Conjunto de perspectivas relevante à questão.\\
\term{CouncilSeat} & Assento sintético alocado a um estrato com justificativa.\\
\term{PerspectiveAgent} & Agente com \term{PerspectiveCharter}, permissões, memória e fontes.\\
\term{DeliberationRound} & Etapa formal com ações permitidas e condições de avanço.\\
\term{ArgumentClaim} & Afirmação passível de sustentação ou contestação.\\
\term{PolicyTension} & Relação argumentada entre objetivos em potencial conflito.\\
\term{ComplianceFinding} & Resultado de norma e evidência por requisito.\\
\term{ScenarioBranch} & Derivação controlada de um cenário com intervenção explícita.\\
\term{SimulationRun} & Execução individual e manifesto congelado.\\
\term{CouncilReport} & Saídas integradas, limites, citações e dissenso.\\\bottomrule
\end{longtable}
\fonte{Elaborado pelos autores (2026).}

\section{Estrutura semântica exemplificativa}
\begin{lstlisting}
SimulationProject -> hasScenario -> ScenarioDefinition
ScenarioDefinition -> usesRegulation -> RegulatorySnapshot
ScenarioDefinition -> evaluates -> SolutionArtifact
ScenarioDefinition -> usesContext -> NationalContextSnapshot
ScenarioDefinition -> instantiatesCouncil -> ShadowCouncil
ShadowCouncil -> hasSeat -> CouncilSeat
CouncilSeat -> assignedAgent -> PerspectiveAgent
PerspectiveAgent -> isGroundedIn -> EvidenceRecord
SimulationRun -> materializes -> WorldSnapshot
SimulationRun -> hasRound -> DeliberationRound
DeliberationRound -> produces -> ArgumentClaim
ArgumentClaim -> supportedBy -> EvidenceRecord
ArgumentClaim -> challenges -> ArgumentClaim
ComplianceFinding -> groundedIn -> LegalProvision
ComplianceFinding -> supportedBy -> EvidenceRecord
CouncilReport -> hasMinorityOpinion -> MinorityOpinion
\end{lstlisting}

\section{Estado jurídico e temporalidade}
O \term{RegulatorySnapshot} deve reter, por provisão, no mínimo: \term{jurisdictionCode}, \term{sourceUri}, \term{legalStatus}, \term{validFrom}, \term{validUntil}, \term{recordedAt}, \term{versionHash}, \term{extractionConfidence}, \term{reviewState} e \term{exceptions}. Valores permitidos para \term{legalStatus}: \term{proposed}, \term{enactedNotEffective}, \term{effective}, \term{suspended}, \term{repealed}, \term{superseded}, \term{unknown}.

\alerta{Um projeto legislativo não pode ser convertido em obrigação vigente. Exemplo de verificação de fonte: na data da especificação, o \href{https://www.camara.leg.br/proposicoesWeb/fichadetramitacao?idProposicao=2487262}{Projeto de Lei (PL) 2338/2023} aparece como proposição em tramitação na Câmara \cite{camara-pl2338}. A plataforma deve consultar o estado oficial em cada execução, em vez de manter um valor permanente fixado.}

\chapter{Composição e representatividade do Shadow Council}
\section{Universo amostral e estratos}
O MVP usa \term{PerspectiveCoverage} com \term{SamplingFrame} de publicações, consultas, estudos, observações e manifestações documentadas. Não usa amostra estatística de seres humanos. \term{CouncilSeat} representa um \emph{recorte de perspectiva} e não uma pessoa real. A Tabela~\ref{tab:estratos} traz uma distribuição ilustrativa de assentos para o Brasil.

\begin{table}[!htb]
\caption{Estratos ilustrativos do Shadow Council para o Brasil}\label{tab:estratos}
\centering
\begin{tabularx}{\textwidth}{@{}L r L@{}}
\toprule\textbf{Estrato ilustrativo (Brasil)} & \textbf{Assentos} & \textbf{Contraste documentado}\\\midrule
Reguladores / fiscalização & 3 & exigibilidade, recursos, poderes e independência\\
Poder público operador & 2 & serviço, compras e responsabilidade administrativa\\
Empresas / startups & 3 & portes, fornecedores, integração e concorrência\\
Pesquisa interdisciplinar & 3 & direito, computação, economia e ciências sociais\\
Sociedade civil e direitos & 3 & direitos digitais, consumidor e acessibilidade\\
Trabalhadores / profissões & 2 & relações de trabalho, supervisão e autonomia\\
Grupos diretamente afetados & 2 & risco vivenciado e acesso a contestação\\
Infraestrutura / ambiente & 2 & energia, dados, soberania e externalidades\\\midrule
\textbf{Total sintético} & \textbf{20} & cotas exploratórias, nunca pesos populacionais\\\bottomrule
\end{tabularx}
\fonte{Elaborado pelos autores (2026). Exemplo sintético.}
\end{table}

\section{Algoritmo de composição}
\begin{enumerate}
\item Construir \term{StakeholderMap}: afetados, operadores, reguladores, titulares de direitos, financiadores e peritos.
\item Definir \term{CoverageDimensions}: papel institucional, exposição ao impacto, porte, território e diversidade de teses.
\item Extrair \term{DocumentedPerspective} de cada estrato; exigir pelo menos uma fonte verificável por posição substantiva.
\item Marcar \term{Unrepresented} para perspectiva sem lastro, sem inventar o que o grupo pensa.
\item Alocar \term{CouncilSeat}, declarando a estratégia de quotas, fontes e lacunas.
\item Gerar \term{PerspectiveCharter} e auditar contradições e conflitos de interesse.
\item Executar \term{BalanceCheck} para identificar cobertura pobre ou dominação de um único estrato/modelo.
\end{enumerate}
Um indicador descritivo possível é \(C = \frac{\text{perspectivas documentadas cobertas}}{\text{perspectivas documentadas consideradas relevantes}}\). O denominador deve ser publicado e justificado. Isso \textbf{não mede} representatividade demográfica nem adesão social.

\section{Persona, limitações e fontes}
Cada \term{PerspectiveCharter} conterá: \term{role}, \term{documentedPositions}, \term{sourceIds}, \term{interests}, \term{constraints}, \term{blindSpots}, \term{allowedActions}, \term{disclosure}, \term{modelVersion}, \term{conflictOfInterest}. O agente deve sinalizar \term{insufficientGrounding} quando uma tese extrapolar as fontes.

\chapter{Mecânica do mundo simulado}
\section{Estado inicial e memórias}
O \term{WorldSnapshot} congela regras, documentos, contexto, atores e dependências. Cada agente possui:\par
\begin{itemize}
\item \term{ReferenceMemory}: posições documentadas e base normativa, \emph{read-only} na execução.
\item \term{WorkingMemory}: objetivos e hipóteses da rodada corrente, descartável.
\item \term{EpisodicMemory}: intervenções, objeções e revisões com origem.
\item \term{ClaimMemory}: grafo de argumentos aceitos, contestados, retirados e condicionais.
\end{itemize}
Memórias privadas de sessões distintas são isoladas. Conhecimentos publicados entre sessões passam por aprovação e revisão de fonte.

\section{Protocolo deliberativo (estado de máquina)}
A Tabela~\ref{tab:protocolo} descreve as etapas do protocolo deliberativo.

\begin{longtable}{@{}p{.14\textwidth}p{.3\textwidth}p{.46\textwidth}@{}}
\caption{Etapas do protocolo deliberativo}\label{tab:protocolo}\\
\toprule\textbf{Etapa} & \textbf{Nome} & \textbf{Transição / saída}\\\midrule\endfirsthead
\toprule\textbf{Etapa} & \textbf{Nome} & \textbf{Transição / saída}\\\midrule\endhead
R0 & \term{Prepare} & Normaliza sementes, congela fonte e verifica lacunas.\\
R1 & \term{IndependentReview} & Cada agente produz posição isolada, sem ler os demais.\\
R2 & \term{OpeningStatements} & Publica teses fundamentadas, ressalvas e perguntas.\\
R3 & \term{CrossExamination} & Questiona afirmações, requer evidências e aponta conflitos.\\
R4 & \term{EvidenceRevision} & Agentes mantêm, revisam ou retiram claims com justificativa.\\
R5 & \term{Negotiation} & Propõem medidas, salvaguardas, exceções e alternativas.\\
R6 & \term{CounterfactualBranches} & Testam alterações controladas sem modificar o estado-base.\\
R7 & \term{MinorityAndAudit} & Preserva dissenso, verifica fontes e documenta limitações.\\
R8 & \term{Publish} & Compila saídas separadas para conformidade e política.\\\bottomrule
\end{longtable}
\fonte{Elaborado pelos autores (2026).}

\noindent\textbf{Interpretação do parâmetro rounds:} a configuração de seis rounds corresponde a R1--R6; R0, R7 e R8 são etapas de preparação, auditoria e publicação fora dessa contagem.\par
\noindent\textbf{Ações permitidas:} \term{assertClaim}, \term{support}, \term{challenge}, \term{requestEvidence}, \term{revise}, \term{withdraw}, \term{proposeAmendment}, \term{declareUncertainty}, \term{submitMinorityOpinion}, \term{abstain}. Cada ação possui \term{actorId}, \term{timestamp}, \term{sourceIds} e \term{roundId}.

\section{Intervenções e ramos contrafactuais}
Exemplo: na hipótese de uma nova regra, o usuário altera \term{complianceCostMultiplier}, \term{inspectionCapacity} ou o texto de uma salvaguarda. A intervenção cria \term{ScenarioBranch} com \term{parentRunId}, \term{patch}, \term{hypothesis} e \term{assumptions}; o estado-base permanece imutável.
Resultados são \emph{efeitos simulados sob suposições}, não previsões do mercado ou da sociedade. A plataforma deve mostrar de que fatores o resultado depende.

\section{Isolamento e repetibilidade}
\begin{itemize}
\item Executar \term{N>=3} repetições em testes exploratórios, com \term{randomSeed} e manifesto por execução (parâmetro proposto, não requisito de fornecedor).
\item Registrar versão e parâmetros dos modelos; arquivar prompts, turnos e respostas para \emph{replay}.
\item Declarar que seeds não asseguram determinismo completo em modelos externos ou chamadas distribuídas.
\item Aplicar limites de custo, timeout, máximo de turnos, budget de tokens e condições de parada por ausência de novos argumentos.
\item Jamais reinterpretar conversa gerada como dado observado de usuários reais.
\end{itemize}

\chapter{Motor jurídico separado do conselho}
\section{Pipeline de conformidade}
\begin{enumerate}
\item \term{ScopeResolution}: determinar jurisdição, papel do provedor/operador, setor, sujeitos afetados e data.
\item \term{ProvisionApplicability}: vincular dispositivos, condições, exceções e eficácia temporal.
\item \term{RequirementDecomposition}: transformar obrigações em controles e testes declarativos com origem.
\item \term{EvidenceMapping}: correlacionar especificação e artefatos com cada requisito.
\item \term{DeterministicChecks}: Web Ontology Language (OWL), Shapes Constraint Language (SHACL) \cite{w3c-shacl} e regras executáveis produzem achados limitados ao que os dados provam.
\item \term{DisputedInterpretation}: encaminhar ambiguidades e conflitos a parecer/revisão.
\item \term{CouncilChallenge}: conselho examina premissas e consequências, sem mudar status jurídico sozinho.
\end{enumerate}
\section{Estados fechados de ComplianceFinding}
A Tabela~\ref{tab:estados} define o vocabulário fechado de estados de \term{ComplianceFinding}.

\begin{table}[!htb]
\caption{Estados fechados de ComplianceFinding}\label{tab:estados}
\centering
\begin{tabularx}{\textwidth}{@{}p{.34\textwidth}L@{}}
\toprule\textbf{Status} & \textbf{Significado}\\\midrule
\term{compliantEvidenced} & Evidência suficiente foi examinada para o controle delimitado (não certificação total).\\
\term{noncompliantEvidenced} & Contrariedade material constatada e fundamentada no requisito.\\
\term{insufficientEvidence} & Não há base para concluir conformidade ou não conformidade.\\
\term{notApplicable} & Regra não incide conforme condições e fontes verificadas.\\
\term{requiresLegalReview} & Ambiguidade, conflito ou exceção exige interpretação humana.\\
\term{futureObligation} & Dispositivo aprovado, mas ainda não exigível na data avaliada.\\
\term{proposalOnly} & Instrumento em discussão, usado apenas como cenário prospectivo.\\\bottomrule
\end{tabularx}
\fonte{Elaborado pelos autores (2026).}
\end{table}

A ausência de um documento não prova violação; \term{noncompliantEvidenced} exige prova positiva do conflito entre fato e obrigação, com aplicabilidade confirmada.

\section{Exemplo de regra semântica (ilustrativa)}
\begin{lstlisting}
# Precondition: Requirement and Assessment exist in validated KG.
# This indicates review of evidence, not a legal finding.
hrai:ComplianceAssessment(?assessment) ^
hrai:hasApplicableRequirement(?assessment, ?requirement) ^
hrai:hasUnverifiedEvidence(?requirement, true)
  -> hrai:requiresEvidenceReview(?assessment, true)
\end{lstlisting}
Para violações por ausência de propriedade (fechamento do mundo), preferir SHACL; para derrogação, prioridades e exceções jurídicas, usar \term{NormativeRuleEngine} com política de resolução explícita, não inferências Semantic Web Rule Language (SWRL) isoladas.

\chapter{Entradas e contratos de dados}
\section{Estrutura do ScenarioRequest}
O arquivo \term{scenario.schema.json} incluído no pacote é o contrato verificável \term{JSON Schema 2020-12}. Os campos fundamentais estão na Tabela~\ref{tab:campos}.

\begin{table}[!htb]
\caption{Campos fundamentais do ScenarioRequest}\label{tab:campos}
\centering
\begin{tabularx}{\textwidth}{@{}p{.36\textwidth}L@{}}
\toprule\textbf{Campo canônico} & \textbf{Uso}\\\midrule
\term{scenarioId} & Identificador estável do cenário.\\
\term{mode} & \term{regulatoryPolicy} ou \term{solutionCompliance}.\\
\term{jurisdiction} & Código de país ISO 3166-1 alpha-2 e nome legível.\\
\term{asOfDate} & Data que determina os requisitos temporais.\\
\term{question} & Pergunta e objetivo analítico.\\
\term{regulatorySources} & Referências normativas, status e versões.\\
\term{contextSources} & Indicadores e relatórios da nação.\\
\term{solution} & Objeto técnico a analisar no modo de solução.\\
\term{council} & Estratos, quotas e parâmetros da seleção.\\
\term{simulation} & Rodadas, repetições, seed e branch interventions.\\\bottomrule
\end{tabularx}
\fonte{Elaborado pelos autores (2026).}
\end{table}

\noindent\textbf{Gate de validação:} no modo \term{solutionCompliance}, \term{solution} e ao menos uma fonte normativa são obrigatórios; se for impossível verificar a fonte ou sua vigência, a execução pode continuar como exploração, mas o relatório não pode declarar conclusão jurídica.

\section{Exemplo mínimo: política regulatória}
\begin{lstlisting}
{
  "scenarioId": "BR-POLICY-001",
  "mode": "regulatoryPolicy",
  "jurisdiction": {"code": "BR", "name": "Brazil"},
  "asOfDate": "2026-11-11",
  "question": "What trade-offs deserve independent review?",
  "regulatorySources": [
    {"id": "PL2338", "title": "PL 2338/2023",
     "status": "proposed",
     "uri": "https://www.camara.leg.br/proposicoesWeb/fichadetramitacao?idProposicao=2487262"}
  ],
  "council": {"strategy": "documentedPerspectiveCoverage",
              "targetSeats": 20},
  "simulation": {"rounds": 6, "replicates": 3,
                 "randomSeed": 42}
}
\end{lstlisting}
\section{Exemplo mínimo: solução e conformidade}
\begin{lstlisting}
{
  "scenarioId": "BR-SOLUTION-001",
  "mode": "solutionCompliance",
  "jurisdiction": {"code": "BR", "name": "Brazil"},
  "asOfDate": "2026-11-11",
  "question": "Which obligations and evidence gaps apply?",
  "regulatorySources": [
    {"id": "LGPD", "title": "Brazilian Data Protection Law",
     "status": "effective", "uri": "https://www.planalto.gov.br/ccivil_03/_ato2015-2018/2018/lei/l13709.htm"}
  ],
  "solution": {"id": "demo-recruitment-01",
               "title": "Recruitment Screening Assistant",
               "summary": "Synthetic recruitment triage specification",
               "artifactType": "briefing"},
  "council": {"strategy": "documentedPerspectiveCoverage",
              "targetSeats": 20},
  "simulation": {"rounds": 6, "replicates": 3,
                 "randomSeed": 17}
}
\end{lstlisting}
Os exemplos são sintéticos. O estado e conteúdo de requisitos legais não são determinados apenas pelo título do documento; devem ser extraídos e conferidos.

\chapter{Saídas, trilha e relatórios}
Cada rodada emite \term{DeliberationEvent} com \term{eventId}, \term{agentId}, \term{roundId}, \term{action}, \term{claimId}, \term{evidenceIds}, \term{timestamp}, \term{contentHash}. Relatórios usam o contrato \term{report.schema.json}. A Tabela~\ref{tab:artefatos} lista os artefatos emitidos por execução.

\begin{longtable}{@{}p{.35\textwidth}p{.55\textwidth}@{}}
\caption{Artefatos emitidos por execução}\label{tab:artefatos}\\
\toprule\textbf{Artefato} & \textbf{Conteúdo}\\\midrule\endfirsthead
\toprule\textbf{Artefato} & \textbf{Conteúdo}\\\midrule\endhead
\term{run-manifest.json} & Versões, sementes, parâmetros, hashes, origem e custo de simulação.\\
\term{world-snapshot.jsonld} & Grafo semântico congelado, entidades, claims, fontes.\\
\term{council-roster.json} & Assentos, estratos, fonte das perspectivas e lacunas de representação.\\
\term{deliberation-events.jsonl} & Todas as ações, turnos, revisões e intervenções.\\
\term{compliance-matrix.json} & Aplicabilidade e status por requisito, com fonte e evidência.\\
\term{tensions.json} & Tensões, argumentos, contrargumentos, evidências e condições.\\
\term{minority-report.md} & Posições minoritárias fundamentadas e não resolvidas.\\
\term{sensitivity-report.json} & Variação de resultados por modelos, composição e seeds.\\
\term{final-report.pdf} & Parecer de simulação, delimitando fatos, hipóteses e limitações.\\\bottomrule
\end{longtable}
\fonte{Elaborado pelos autores (2026).}

\noindent\textbf{Destaque de saída:} um sumário \term{CouncilSupport} (ex.: 12/20 agentes sintéticos defendem uma tese) pode ser exibido somente como contagem da execução. \textbf{Nunca} como percentual de cidadãos, eleitores ou stakeholders reais.

\chapter{API e eventos de execução}
A Tabela~\ref{tab:api} lista os endpoints propostos da interface de programação de aplicações (API, \textit{Application Programming Interface}), no estilo Representational State Transfer (REST).

\begin{table}[!htb]
\caption{Endpoints REST propostos e seu comportamento}\label{tab:api}
\centering
\begin{tabularx}{\textwidth}{@{}p{.42\textwidth}L@{}}
\toprule\textbf{Endpoint REST (proposta)} & \textbf{Comportamento}\\\midrule
\term{POST /v1/scenarios} & Criar cenário e validar contrato.\\
\term{POST /v1/scenarios/:id/ingest} & Ingerir e verificar fontes e briefings.\\
\term{POST /v1/scenarios/:id/preview} & Mostrar contexto, escopo e lacunas antes da execução.\\
\term{POST /v1/scenarios/:id/council} & Construir e auditar conselho.\\
\term{POST /v1/scenarios/:id/runs} & Iniciar execução com manifesto imutável.\\
\term{GET /v1/runs/:id/events} & Consultar stream de eventos com cursor.\\
\term{POST /v1/runs/:id/branches} & Criar contrafactual derivado sem alterar original.\\
\term{GET /v1/runs/:id/findings} & Obter achados e matrizes.\\
\term{GET /v1/runs/:id/reports} & Obter versões JSON/MD/PDF.\\
\term{POST /v1/runs/:id/stop} & Interromper a execução com auditoria.\\\bottomrule
\end{tabularx}
\fonte{Elaborado pelos autores (2026).}
\end{table}

A fila de eventos pode usar \term{scenario.ingested}, \term{snapshot.frozen}, \term{council.prepared}, \term{round.started}, \term{claim.challenged}, \term{round.completed}, \term{run.completed}, \term{report.published}. Executores devem tratar idempotência, reentrega, paralelismo e ordenação por \term{runId}.

\chapter{Interface de usuário (SCSE Studio)}
A interface precisa permitir: (1) seleção do modo; (2) upload de documentos e escolha de jurisdição/data; (3) pré-visualização do \term{RegulatorySnapshot}; (4) matriz de assentos e fontes por estrato; (5) linha do tempo das rodadas; (6) grafo de argumentos com apoios/contestações; (7) \emph{compare branches}; (8) inspeção de cada \term{ComplianceFinding}; (9) exportação; (10) questionamento de agentes com rastreio do contexto.

\destaque{\textbf{Princípio de experiência do usuário (UX).} Toda resposta deve distinguir visualmente \term{VerifiedLegalFact}, \term{EvidenceBasedFinding}, \term{SyntheticPerspective}, \term{Hypothesis} e \term{Unknown}. Uma interface que misture essas categorias perderia justamente a integridade que o produto pretende oferecer.}

\chapter{Segurança, direitos e governança}
\begin{itemize}
\item \textbf{Acesso:} multi-tenant, escopos \term{project:read}, \term{run:execute}, \term{source:review}, \term{agent:register}, \term{report:publish}; controle de acesso baseado em papéis e em atributos (RBAC/ABAC, \textit{role-based} e \textit{attribute-based access control}) por tipo de evidência.
\item \textbf{Dados:} cifragem, retenção por classe, mascaramento de informações pessoais e isolamento de fontes privadas; proibir publicação inadvertida de dossiês.
\item \textbf{Prompt injection:} tratar textos externos como dados não confiáveis; ferramentas permitidas por agente e política de execução.
\item \textbf{Conflito de interesse:} registrar patrocínio, afiliações autorizadas e potenciais dependências de fornecedor.
\item \textbf{Supervisão humana:} interpretações jurídicas materiais, publicação oficial e intervenções normativas exigem validação humana.
\item \textbf{Não impersonation:} agentes não se passarão por pessoas, grupos ou reguladores reais; perspectivas sintéticas serão rotuladas como tal.
\item \textbf{Auditoria:} logs encadeados por hash, controle de versões e rastreio de intervenções do orquestrador.
\end{itemize}

\chapter{Requisitos verificáveis e aceite}
A Tabela~\ref{tab:requisitos} lista os requisitos funcionais, seus critérios de aceite e prioridades.

\begin{longtable}{@{}p{.11\textwidth}p{.69\textwidth}p{.1\textwidth}@{}}
\caption{Requisitos funcionais e critérios de aceite}\label{tab:requisitos}\\
\toprule\textbf{ID} & \textbf{Critério de aceite} & \textbf{Prior.}\\\midrule\endfirsthead
\toprule\textbf{ID} & \textbf{Critério de aceite} & \textbf{Prior.}\\\midrule\endhead
FR-01 & Receber os dois modos com JSON Schema e relatar erros por campo. & P0\\
FR-02 & Resolver fonte, jurisdição, versão, status e data antes de liberar achados jurídicos. & P0\\
FR-03 & Congelar \term{WorldSnapshot} e manifesto identificados por hash. & P0\\
FR-04 & Produzir \term{CouncilRoster} por estratos, fontes e lacunas de cobertura. & P0\\
FR-05 & Executar pelo menos seis etapas com eventos auditáveis de discussão e revisão. & P0\\
FR-06 & Preservar opiniões minoritárias não resolvidas no relatório. & P0\\
FR-07 & Emitir estados de conformidade apenas do vocabulário controlado. & P0\\
FR-08 & Marcar normas \term{proposed} como \term{proposalOnly}. & P0\\
FR-09 & Gerar matriz de achados por requisito, fontes e evidências associadas. & P0\\
FR-10 & Gerar relatório com declaração de natureza sintética e limitações. & P0\\
FR-11 & Suportar \term{ScenarioBranch} com diffs e estado-base imutável. & P1\\
FR-12 & Reexecutar cenário em N réplicas e medir sensibilidade. & P1\\
FR-13 & Permitir comparar dois regimes normativos sem afirmar equivalência jurídica automática. & P1\\
FR-14 & Integrar agentes externos por manifesto de capacidades e permissões. & P1\\
FR-15 & Exportar \term{json}, \term{md}, \term{pdf} e logs de evento. & P0\\
FR-16 & Reportar \term{unknown} quando não houver fonte ou documento necessário. & P0\\\bottomrule
\end{longtable}
\fonte{Elaborado pelos autores (2026).}

\chapter{Avaliação experimental e qualidade}
\section{Metodologia proposta}
Testar três condições sobre casos de referência previamente auditados: \textbf{B1} agente único, \textbf{B2} debate de agentes sem estratificação e \textbf{SCSE} conselho estratificado e rastreável. Especialistas independentes devem construir previamente rubricas e um conjunto de achados verificáveis, mantidos cegos para os modelos quando possível. A Tabela~\ref{tab:metricas} define as métricas e sua operacionalização.

\begin{table}[!htb]
\caption{Métricas de avaliação e operacionalização}\label{tab:metricas}
\centering
\begin{tabularx}{\textwidth}{@{}p{.33\textwidth}L@{}}
\toprule\textbf{Métrica} & \textbf{Operacionalização}\\\midrule
\term{ClaimEvidencePrecision} & Claims relevantes corretamente sustentados / claims auditados.\\
\term{RelevantGapRecall} & Lacunas de referência detectadas / lacunas de referência.\\
\term{FalseLegalFindingRate} & Achados jurídicos falsos ou injustificados / achados emitidos.\\
\term{PerspectiveCoverage} & Cobertura dos eixos/estratos do frame documentado.\\
\term{DissentRetention} & Teses minoritárias relevantes preservadas / teses minoritárias identificadas.\\
\term{RunSensitivity} & Dispersão entre repetições, prompts, modelos e ordem de falas.\\
\term{ExpertUtilityScore} & Utilidade do relatório julgada por especialistas com rubrica publicada.\\
\term{CostPerValidatedFinding} & Custo total / achados validados independentemente.\\\bottomrule
\end{tabularx}
\fonte{Elaborado pelos autores (2026).}
\end{table}

\section{Limites científicos}
O conselho não estima distribuição real de preferências; modelos compartilham vieses, corpora e pressupostos, podendo produzir falso consenso. Escores de apoio dentro da simulação são apenas resultados internos. Para inferir comportamento real é necessária validação com pessoas e métodos empíricos adequados.

\chapter{MVP e ordem de implementação}
A Tabela~\ref{tab:fases} apresenta as fases de implementação, suas entregas e gates.

\begin{longtable}{@{}p{.16\textwidth}p{.38\textwidth}p{.35\textwidth}@{}}
\caption{Fases do MVP, entregas e gates}\label{tab:fases}\\
\toprule\textbf{Fase} & \textbf{Entrega} & \textbf{Gate}\\\midrule\endfirsthead
\toprule\textbf{Fase} & \textbf{Entrega} & \textbf{Gate}\\\midrule\endhead
M0 & Schemas, modos, armazenamento, ingestão, versionamento e snapshots. & Fontes e timestamps auditáveis.\\
M1 & Conselho de 20 assentos documentais, personas, 6 rounds, logs, relatório minoritário. & Cobertura e integridade avaliadas por revisão humana.\\
M2 & Motor de matriz de conformidade, SHACL, estados e explicações por requisito. & Nenhum achado deriva apenas de voto sintético.\\
M3 & Branches, multi-run, sensibilidade, agentes externos e interface visual. & Reprodução e isolamento entre projetos demonstrados.\\\bottomrule
\end{longtable}
\fonte{Elaborado pelos autores (2026).}

\noindent\textbf{MVP de ponta a ponta:} carregar fonte legislativa brasileira em \term{proposed}, dossiê de solução fictícia e normas efetivas selecionadas; compor conselho; rodar debate; emitir matriz de requisitos com estado \term{insufficientEvidence}/\term{proposalOnly}, objeções, relatório de minoria, fontes e manifesto.

\chapter{Cenários pedagógicos de ponta a ponta}
\section{Cenário 1: revisão de política no Brasil}
\textbf{Seed:} texto de proposta regulatória, seu status oficial, Plano Brasileiro de Inteligência Artificial (PBIA) \cite{pbia2025}, fontes econômicas e documentos públicos de stakeholders. \textbf{Pergunta:} ``Quais consequências da obrigatoriedade de auditoria ex-ante merecem investigação?'' Agentes de direitos digitais e trabalhadores podem contestar uma redução de salvaguardas; representantes sintéticos de pequenas empresas podem questionar custos; peritos em fiscalização discutem capacidade institucional. A saída é uma matriz de tensões fundamentadas, sem afirmar que a maioria dos brasileiros apoia algum regime.

\section{Cenário 2: briefing de recrutamento automatizado}
\textbf{Seed:} \term{RecruitmentScreeningAssistant} (fictício), com dados de candidatos, triagem e explicações ainda não documentadas; marco de proteção de dados brasileiro selecionado e datado \cite{lgpd2018}. \textbf{Pergunta:} ``Que evidências são necessárias e quais riscos devem ser analisados?'' O motor legal identifica requisitos candidatos e demanda validação de aplicabilidade; o conselho investiga riscos de viés, explicabilidade, direitos dos titulares, efeitos trabalhistas e viabilidade de mitigação. O relatório não declarará violação se apenas faltar documentação.

\section{O que deve ser entregue numa execução}
\begin{enumerate}
\item \term{ScopeStatement}: modo, norma, país, data e questão.
\item \term{SourceInventory}: fontes verificadas, insuficientes e propostas.
\item \term{CouncilRoster}: estratos, fontes e lacunas.
\item \term{ArgumentGraph}: teses, objeções e revisões com citações.
\item \term{ComplianceMatrix} (somente modo de solução).
\item \term{PolicyTensionMatrix} (modo regulatório e quando pertinente ao modo de solução).
\item \term{MinorityReport}, \term{Unknowns}, \term{Recommendations} condicionais.
\item \term{RunManifest} e disponibilidade de inspeção por evento.
\end{enumerate}

\chapter{Considerações finais}
\destaque{\textbf{Status desta entrega:} especificação arquitetural e funcional, não uma implementação executada. Os JSON Schemas e exemplos anexados são contratos de engenharia propostos. Toda conclusão legal dependerá de fontes oficiais atualizadas, de uma regra aplicável e da evidência concreta do caso.}

\section{Fontes e referências de implementação}
\begin{enumerate}
\item \textbf{Fonte técnica:} MiroFish, repositório original \cite{mirofish}. Fluxo de graph building, personas, simulação, relatórios e deep interaction.
\item \textbf{Fonte técnica:} OASIS, da CAMEL-AI \cite{camel-oasis}. Motor aberto de simulação social.
\item \textbf{Fonte oficial:} Câmara dos Deputados, tramitação do PL 2338/2023 \cite{camara-pl2338}.
\item \textbf{Fonte normativa de exemplo:} Lei Geral de Proteção de Dados Pessoais (LGPD), Lei 13.709/2018 \cite{lgpd2018}.
\item \textbf{Documento de referência fornecido:} \emph{Guia de Conformidade e Governança para Sistemas de Inteligência Artificial de Alto Risco (HRAI)} \cite{estudo-guia-hrai}.
\item \textbf{Documento de referência fornecido:} \emph{Análise Comparativa e Diretrizes Estratégicas para a Governança e Plano de Inteligência Artificial no Brasil} \cite{estudo-comparativo}.
\item \textbf{Documento de referência fornecido:} \emph{Brainstorming de Artefatos Técnicos Atemporais para Governança de IA} \cite{estudo-brainstorm}.
\item \textbf{Alinhamento interno:} \emph{HRAI-GO --- Plataforma Aberta de Inteligência Regulatória e Governança de IA}, documento pedagógico e especificação preliminar do Shadow Council, 11 nov. 2026 (DC-PLAT-01, \texttt{docs/hrai-go/}).
\end{enumerate}

% ============================================================ ELEMENTOS PÓS-TEXTUAIS
\postextual
\bibliography{referencias,refs-local}

\end{document}
